\hspace{3mm}

\centerline{\bf \Huge  Задачи по графам.}
\centerline{\bf \Huge Уровень 3:}
\centerline{\bf \Huge Калмыцкие Ханы и Дагестанские Горцы}

\begin{flushright}
    \scriptsize{В Дагестане мы очень любим калмыцкий чай.\\
                Называем мы его "калмух чай".}
\end{flushright}

\problem{1} 
%$\bigl[\textbf{Курчатов, 2017, 8 класс}\bigr]$
В математическом конкурсе участвуют 14 школьников. Участникам конкурса было предложено 6 задач. В результате каждую задачу правильно решили больше половины школьников. Докажите, что обязательно найдётся пара участников, которые в объединении правильно решили все задачи.

\problem{2} 
%\bigl[\textbf{Высшая Проба, 2017, 9 класс}\bigr]$
Каждый член партии доверяет пяти однопартийцам, но никакие двое не доверяют друг другу. При каком минимальном размере партии такое возможно?

\problem{3} 
%$\bigl[\textbf{Олимпиада СПбГУ, 2022, 10-11 класс. Отборочный этап.}\bigr]$ 
На доске написаны 1000 различных натуральных чисел. Оказалось, что для каждого написанного числа a на доске найдется еще хотя бы одно число b такое, что |$a - b| - $ простое число. Докажите, что можно подчеркнуть не более 500 чисел так, чтобы для каждого неподчеркнутого числа a нашлось подчеркнутое число b, для которого |$a - b| -$ простое число.

\problem{4}
%$\bigl[\textbf{Турнир городов, 2009, сложный тур, 8-9}\bigr]$
Оля и Максим оплатили путешествие по архипелагу из 2009 островов, где некоторые острова связаны двусторонними маршрутами катера. Они путешествуют, играя. Сначала Оля выбирает остров, на который они прилетают. Затем они путешествуют вместе на катерах, по очереди выбирая остров, на котором еще не были (первый раз выбирает Максим). Кто не сможет выбрать остров, проиграл. Докажите, что Оля может выиграть.

\problem{5}
%$\bigl[\textbf{Всеросс, 1993, округ, 11.8}\bigr]$
В стране 1993 города, и из каждого выходит не менее 93 дорог. Известно, что из каждого города можно проехать по дорогам в любой другой. Докажите, что это можно сделать не более, чем с 62 пересадками. (Дорога соединяет между собой два города.)